%% ------------------------------------------------------------------
%% Problema 7
%% ------------------------------------------------------------------
\section{Problema 7. Descuento en medicamentos}

% TODO: modelar la determinaci\'on del descuento en la compra de medicamentos
% con grafos de dependencia, COBEGIN--COEND y condiciones de Bernstein.
\documentclass[12pt]{article}
\usepackage[utf8]{inputenc}
\usepackage[spanish]{babel}
\usepackage{amsmath, amssymb, amsthm}
\usepackage{tikz}
\usetikzlibrary{arrows.meta, positioning, shapes.geometric}
\usepackage{geometry}
\geometry{letterpaper, margin=1in}

\begin{document}

\section{Problema 7. Redes de Interconexión y Topologías}

\subsection{Análisis Teórico: Hipercubos de Dimensión $n$}
Un hipercubo $n$-dimensional ($Q_n$) es un grafo cuyos vértices corresponden a las $2^n$ cadenas binarias de longitud $n$. Dos vértices están conectados por una arista si y solo si sus representaciones binarias difieren exactamente en un solo bit.

\subsubsection{Propiedades de la Topología}
\begin{itemize}
    \item \textbf{Número de Vértices (Nodos):} $V = 2^n$
    \item \textbf{Número de Aristas (Enlaces):} $E = n \cdot 2^{n-1}$
    \item \textbf{Grado de cada Nodo:} $D = n$
    \item \textbf{Diámetro de la Red:} $\text{Diámetro} = n$
\end{itemize}

\subsection{Representación Gráfica}
A continuación se ilustra la estructura de un hipercubo tridimensional ($Q_3$), conformado por $2^3 = 8$ nodos y $12$ enlaces binarios.

\begin{figure}[h!]
\centering
\begin{tikzpicture}[scale=2, auto, >=stealth,
    node distance=2cm,
    vertex/.style={circle, draw=blue!80, fill=blue!10, thick, inner sep=2pt, minimum size=8mm}]

    % Nodos del cubo frontal
    \node[vertex] (000) at (0,0) {000};
    \node[vertex] (100) at (2,0) {100};
    \node[vertex] (110) at (2,2) {110};
    \node[vertex] (010) at (0,2) {010};

    % Nodos del cubo trasero
    \node[vertex] (001) at (0.8,0.8) {001};
    \node[vertex] (101) at (2.8,0.8) {101};
    \node[vertex] (111) at (2.8,2.8) {111};
    \node[vertex] (011) at (0.8,2.8) {011};

    % Aristas del cubo frontal
    \draw[thick] (000) -- (100);
    \draw[thick] (100) -- (110);
    \draw[thick] (110) -- (010);
    \draw[thick] (010) -- (000);

    % Aristas del cubo trasero
    \draw[thick] (001) -- (101);
    \draw[thick] (101) -- (111);
    \draw[thick] (111) -- (011);
    \draw[thick] (011) -- (001);

    % Aristas de interconexión (frente a atrás)
    \draw[thick, dashed, red!70!black] (000) -- (001);
    \draw[thick, dashed, red!70!black] (100) -- (101);
    \draw[thick, dashed, red!70!black] (110) -- (111);
    \draw[thick, dashed, red!70!black] (010) -- (011);

\end{tikzpicture}
\caption{Topología de Hipercubo 3D ($Q_3$). Las líneas punteadas rojas unen nodos que difieren en el tercer bit.}
\end{figure}

\subsection{Conclusiones}
La topología en hipercubo ofrece un excelente compromiso entre el diámetro de la red y el grado de los nodos, permitiendo un enrutamiento eficiente de complejidad $\mathcal{O}(n)$ en algoritmos de paso de mensajes paralelos.

\end{document}
